\documentclass[11pt]{article}
\usepackage[margin=0.9in]{geometry}
\usepackage{amsmath,amssymb,amsthm,booktabs,longtable,array}
\usepackage[T1]{fontenc}
\usepackage[utf8]{inputenc}
\usepackage{lmodern,microtype}
\usepackage{xurl}
\usepackage[colorlinks=true,linkcolor=blue,urlcolor=blue]{hyperref}
\providecommand{\hrbpath}[1]{\nolinkurl{#1}}
\providecommand{\hncpath}[1]{\nolinkurl{#1}}
\newcommand{\norm}[1]{\left\lVert#1\right\rVert}
\newtheorem{theorem}{Theorem}[section]
\newtheorem{lemma}[theorem]{Lemma}
\newtheorem{proposition}[theorem]{Proposition}
\newtheorem{corollary}[theorem]{Corollary}
\theoremstyle{remark}
\newtheorem{remark}[theorem]{Remark}
\setlength{\emergencystretch}{5em}
\setcounter{tocdepth}{2}
\author{Research handoff}
\date{17 October 2026}

\title{Four dimensional Yang--Mills\ research dossier}
\begin{document}
\maketitle
\tableofcontents
\clearpage
\section{Purpose, conclusion and evidence rules}
\label{hdc-purpose}

This dossier transfers a sustained attempt to complete a rigorous
existence argument for a nontrivial four-dimensional Yang--Mills
quantum theory. The project starts from the repository
\emph{Yang--Mills-LLM-experiment}, hosted at
\texttt{ethanwsussman/Yang--Mills-LLM-experiment} with a README
note signed LA, and from the Ba{\l}aban renormalization literature. The detailed
baseline assessment below distinguishes the repository's stated
ambition from the precise scope of its displayed results.

At handoff there is a substantial local ultraviolet construction
in the working proof chain, together with new exact statements
about full finite Wilson measures. There is not yet a completed
continuum existence proof. A complete-density multiscale theorem,
observable estimates, and continuum identification still require
mathematical arguments.

The final research step, Notebook 39, completes a restricted
normalized replacement operation on arbitrary collections of
separated fixed boxes in the actual bare Wilson theory. Its
denominator is controlled, its exterior marginal is preserved,
and its effect on local bounded observables does not grow with
the number of distant boxes. This is a natural stopping point
for the present round: one finite-measure operation is complete,
and its remaining distance from a multiscale construction is explicit.

\subsection{What a proof claim in this bundle means}

The documents distinguish five kinds of statement.
\begin{enumerate}
\item A \emph{derived working result} has a written argument in the
included notebooks, with its assumptions and scope stated. Many
recent modules have separate internal mathematical audits.
This is not independent professional verification of all earlier
dependencies.
\item A \emph{conditional reduction} proves an implication from
explicit model bounds. Its hypotheses are not counted as established
merely because the reduction is proved.
\item A \emph{published input} is tied to an actual primary statement.
Its use requires matching its domains, normalizations and uniformities.
\item An \emph{open interface} is a step of this proposed route
for which no complete argument is supplied here. A different
route to existence might replace it.
\item A \emph{superseded argument} is retained for provenance, with
its correction or replacement identified. Its presence in the
research archive is not an endorsement.
\end{enumerate}

The terms our addition and derived here identify contributions
to this project record. They are not claims of novelty relative
to the whole literature. The compact-group Laplace, tree-gauge,
coupling and chessboard mechanisms have standard antecedents.
Their precise assembly and use in the present program are what
this dossier records.

\subsection{The archive's order of authority}

Read the final handoff documents for the current map, then the
numbered notebooks for statements and proofs. Read each
notebook's dependencies before treating its theorem as
unconditional. The research archive contains fragments,
calculation notes, audits and historical checkpoints. Later
corrections override earlier speculation and earlier next-task
paragraphs.

The catalogue supplies an entry for every numbered notebook:
its location, central results, assumptions, dependencies and
limitations. The manifest identifies exact files by SHA-256.
The repository is pinned separately; a future remote update
is not silently substituted for the examined commit.

\section{The target and the hierarchy of conclusions}
\label{hdc-target}

The working regulator is a finite four-dimensional periodic
Wilson lattice gauge theory, principally for \(SU(2)\) in the
inherited construction. At positive coupling its compact group
integral exists automatically. The hard question is a nontrivial
continuum theory at vanishing lattice spacing, with the
quantum-field properties needed to identify it as Yang--Mills.

Three conclusions must be distinguished. A continuum limit of
separated scalar action-density correlations is a useful
observable sector. Reconstruction from a sufficient Euclidean
hierarchy could give a local relativistic quantum observable
theory. Identifying it as nontrivial Yang--Mills requires an
adequate gauge-invariant algebra and its relation to the lattice
theory and short-distance normalization. A physical mass gap is
a further spectral assertion. A result for \(SU(2)\) does not
by itself cover every compact simple gauge group.

Subsequential existence is a legitimate target. A unique cutoff
limit or preferred contact-term convention is not imposed
solely for convenience. What is indispensable is one consistent
limiting hierarchy with the requisite positivity, covariance,
regularity and physical interpretation.

\subsection{Distinct parameters and uniformities}

In the main local chain, \(M\) denotes preceding homogeneous
classical depth; \(\ell\) is the fixed primitive ratio;
\(R=\ell^s\) is a fixed composite quantum ratio; \(J\) is a
fixed auxiliary grouping; \(N\) is a finite period; \(b\) is
a small analytic radius; \(g\) is the running coupling; and
\(K\) is a large endpoint ratio in determinant comparisons.
In physical lattice statements \(\eta\) denotes mesh size.
Individual theorem statements control any changes of notation.

Uniformity in \(M,N\) on a fixed small ball is not uniformity on
arbitrary exterior fields. Constants depending on fixed \(R,J\)
are not uniform as either grows. A fixed-box bound uniform in
surrounding volume is not a bound on a physical box containing
more and more lattice sites. An infinite structural kernel is
not automatically the actual finite winding kernel. Required
old-depth shifts, endpoint ratios and source derivatives must
belong to one joint structural limit.


% Body fragment. Requires amsmath, amssymb and hyperref. Prefix hrb-.
% Repository frozen at e392ef45013a581664bab76b90867170b30fe10e.
\providecommand{\hrbpath}[1]{\texttt{\detokenize{#1}}}

\section{Repository baseline, mathematical provenance, and additions}
\label{hrb-section}

\noindent\textbf{Evidence date: 17 October 2026.}
This report distinguishes what the deposited repository claims;
what its displayed arguments take as inputs; what published
literature supplies in its own setting; and what this research
series adds. It is a reading and dependency report, not
certification of either the repository or our notebooks.
Neither manuscript length, a theorem environment, successful
compilation, nor repeated internal review establishes a proof.

The repository claims an apparent existence argument on
\(\mathbb R^4\). Describing the baseline as only finite-volume
work, only a programme, or no claimed existence argument would
be inaccurate. Identifying that claim with an externally verified
Yang--Mills construction with a mass gap would also be inaccurate.
The shorter manuscript expressly states a narrower observable
theorem, records conditional deductions and excludes stronger
conclusions. Our additions address selected interfaces and
develop a separately specified local construction. They do not
retrospectively verify every argument in the deposit.

\subsection{Frozen evidence and reproducibility}
\label{hrb-version}

The source is
\url{https://github.com/ethanwsussman/Yang--Mills-LLM-experiment},
frozen at \hrbpath{e392ef45013a581664bab76b90867170b30fe10e}.
The commit is dated 16 October 2026, 02:33:55 \(-0400\),
has message \emph{Update README.md}, and records the Git author
as alpoge. This identifies a version, not authorship of every
generated or human-written argument. The raw master has an
empty author field. We do not infer a formal author list from
the hosting account or Git metadata. The checkout was clean
when checked for this report and was not edited. Later remote
changes are outside the evidence boundary.

Raw line numbers are counted from one in that commit.
PDF locators are printed pages of
\hrbpath{ym4_manuscript_zeroth_draft.pdf}, dated 6 October 2026.
In the local extraction its page markers agree with printed
pages. Raw labels are preferable to generated theorem numbers:
the raw master does not presently assemble unchanged.
Its lines 38--40 input
\hrbpath{theorem_0_complete_account_not_done_for_ethan_v2_part1},
\hrbpath{theorem_0_complete_account_not_done_for_ethan_v2_part2},
and \hrbpath{theorem_0_complete_account_not_done_for_ethan_v2_part3},
whereas the actual files are
\hrbpath{theorem_0_complete_account_raw_I.tex},
\hrbpath{theorem_0_complete_account_raw_II.tex}, and
\hrbpath{theorem_0_complete_account_raw_III.tex}.
This is a reproducibility naming issue, not a mathematical
counterexample. The handoff preserves it rather than silently
repairing third-party source.

The three raw parts have respectively 166,711, 145,624 and
111,583 lines: 423,918 in total. These counts measure material
to audit, not proved mathematics. The companion
\hrbpath{handoff_repository_source_map.json} records commit,
SHA-256 hashes, byte counts, exact locators and selected links.
The handoff uses a manifest and retrieval instructions rather
than redistributing third-party source PDFs or articles.

Below, I, II and III mean the three raw files; Draft means
the zeroth PDF; Our Notebook always means a document of
this research series, not a repository document.

\subsection{What the README and zeroth manuscript claim}
\label{hrb-claims}

\paragraph{README, lines 1--6.}
The opening describes ongoing model-assisted QED3, YM3 and
YM4 work, with exposition and formalization in progress.
The note at lines 5--6, signed LA, explicitly reports an
apparent YM4 existence argument in \(\mathbb R^4\), emphasizes
its length and lack of readability, and explains the simpler
presentations as a way of exposing ideas, including a
reflection-positivity route to rotations. This is a qualified
public claim, not a theorem statement or a mass-gap claim.

The pinned tree also contains
\hrbpath{QED3_slop/YM_Rotation_Invariance_simplified.pdf},
\hrbpath{QED3_slop/qed3_manuscript.pdf}, and
\hrbpath{QED3_slop/qed3.pdf}. Their presence and the
README already place the reflection-based rotation
idea in the baseline. This report does not independently
certify those lower-dimensional manuscripts or treat
our pure-Wilson interface work as the origin of that idea.

\paragraph{Draft, abstract and Sections 1.1--1.4, pp.1--23.}
The object is the Wilson lattice measure for \(SU(2)\), on
compatible periodic tori, and correlations of one local
gauge-invariant observable: smeared action density.
Supports of the test functions are pairwise disjoint.
The density has canonical dimension four; its leading
separated two-point kernel has singularity \(|x-y|^{-8}\).
There is no unmentioned prescription extending it across
coincidence.

In the draft's convention,
\[
 A(U)=\sum_p(1-\operatorname{Re}\operatorname{tr}U(\partial p)),
 \qquad d\mu(U)\propto e^{-g_0^{-2}A(U)}\,dU,
 \qquad \operatorname{tr}1=1 .
\]
The observable is the corresponding test-weighted plaquette
sum. Complex sources change action weights; connected
correlations are derivatives of the logarithm of the
sourced partition function, and centred moments are
assembled from cumulants. This convention matters when
comparing beta coefficients and two-point amplitudes.
See Draft pp.7--9 and I, Chapter 1, beginning at line 401.

\paragraph{Theorem 0: scope of the clauses.}
This is a navigation map, not a substitute for the
parameter hypotheses in the displayed theorem.
\begin{description}
\item[(0), running coupling.]
The draft asserts an actual recursively generated trajectory
stopped at a first-passage scale, with two-sided cumulative
inverse-coupling bounds and, for \(SU(2)\), a cumulative
one-loop law. It does not require a positive decrement at
every step. I:2533--2554 explicitly distinguishes cumulative
drift from a single-step sign.
\item[(i), ultraviolet stability.]
The application of Ba{\l}aban's complete-density theorem
uses the weak-coupling containment furnished by (0).
Stability concerns effective densities and their prescribed
representation, not by itself a continuum observable limit.
\item[(iii), sourced correlations.]
Connected action-density bounds are asserted uniformly in
cutoff; a separate per-ball estimate supplies the needed
volume independence. This source theorem and its running
shift are substantive inputs, not consequences of a vacuum
partition-function estimate alone.
\item[(ii), limits and reconstruction.]
The claim includes torus and \(\mathbb R^4\) limits with
the stated fixed-order regularity, translations, hypercubic
symmetry, permutation symmetry and reflection positivity.
It reconstructs a Hilbert space, vacuum and nonnegative
time-translation generator in the stated sense.
The \(\mathbb R^4\) limit is subsequential. Fixed-torus
convergence without subsequences and identification of
first-passage limits use (v), inheriting its dependencies.
\item[(iv), a two-point witness.]
At small separated distances the connected two-point
function is comparable to the specified positive free
curvature kernel times the fourth power of a running
coupling, with controlled relative error. A one-torus
witness is transferred to other volumes. Its consequences
include a nonzero centred reflected vector and logarithmic
short-distance behavior.
\item[(v), a torus family.]
Under comparison inputs, a unique analytic one-parameter
family is asserted at fixed torus size. Two-cutoff and
first-passage/tuned identifications are stronger than
subsequential compactness. They are not infinite-volume
uniqueness or clustering.
\item[(vi), rotations.]
Full \(O(4)\) invariance is conditional on a displayed
far-sphere angular-momentum-flux hypothesis, which the
draft does not establish. Short-distance approximate
rotation statements and the README's later reflection
argument are not interchangeable assertions.
\end{description}

Exact reading locations are Draft Theorem 0, pp.12--23;
Sections 2, 3, 4, 5--6, 7 and 8 for (0)--(i), (iii),
(ii), (iv), (v) and (vi), respectively. The raw main
statement is I:2512, \hrbpath{3.1:thm-main}; its dependency
assembly is III:110784--110860, \hrbpath{ch:11}.

\paragraph{The meaning of non-triviality is deliberately limited.}
Draft p.5 and I:2475--2510,
\hrbpath{3.1:def-non-trivial}, distinguish the positive,
logarithmically running two-point law from stronger
interaction assertions. The raw definition expressly
permits a generalized-free realization of an admissible
two-point function and observes that a composite of a free
field need not be Gaussian. Our counterexamples and
spectral-witness discussion must not be presented as
discovering a limitation that the baseline denied.
They clarify what a stronger Yang--Mills identification
would require.

\paragraph{Explicit exclusions.}
Draft pp.5--6 and 150--151 do not claim a mass gap,
clustering, uniqueness of the infinite-volume limit or
vacuum, coincidence extensions, Wightman fields with the
required growth in the number of arguments, Wilson-loop
or block-observable laws, unconditional full rotations,
or the complete result for \(SU(N)\), \(N\geq3\).
Its small-torus spectral discussion and finite-size
criterion in the outlook are not a mass gap on
\(\mathbb R^4\). The title's construction and
non-triviality language has these express limitations.

\subsection{Proof ledger and unresolved version crosswalk}
\label{hrb-ledger}

The hypotheses (H1)--(H7) in Draft Section 1.3 organize
geometry, thresholds and quantifier order. They are not
all the dynamical inputs. Section 9, pp.141--151, is the
additional ledger: published citations, statements used
as Ba{\l}aban asserts them, companion proofs, deferred
proofs and unfinished supporting statements are separated.

\paragraph{What the shorter draft says is written.}
Draft pp.6 and 143 says that no own input below \(SU(2)\)
clauses (0)--(iv) is deferred, except portions of (ii)
using (v). It claims companion proofs for the source
theorem, witness decomposition, Gaussian extraction and
one-cell large-field estimates, among others.
Section 9.1 nevertheless records fifteen imported
statements/readings without a printed proof covering
their asserted use. A written companion proof and
independent verification of every imported generality
are different propositions.

\paragraph{What particular raw sites say.}
III:110862--110870 explicitly says no clause is proved
outright. III:111167--111179 organizes outstanding proofs
into five categories. This concerns implications from
named inputs; it does not mean no part contains a proof.
Two important exact sites are:
\begin{enumerate}
\item I:3901--3945, \hrbpath{4.2:lem-one-loop}, states
the finite Gaussian one-loop law by citation to
CMP109/CMP116 and begins an outline:
\(\left|\sum_{l=1}^k\beta_l^0-kc_0\right|\leq c_2\).
I:3895--3898 gives \(c_0=11N^2(12\pi^2)^{-1}\log L\)
in the deposited normalized-trace convention.
\item I:4031--4088, \hrbpath{4.2:lem-unmarked}, marks
the CMP116 remainder transfer to the CMP119 cutoff
and fluctuation radius as an incomplete step.
I:4136--4198, \hrbpath{4.3:thm-small-field-step},
displays the stronger family-form/cutoff reading.
I:56419--56440, \hrbpath{4.19:prop-clause-zero},
lists the relevant assumptions; lines 56428--56430
call both one-loop and unmarked statements incomplete,
although the current one-loop lemma instead uses
the citation/outline description above.
\end{enumerate}

The shorter draft's completion claim and these raw
markers are not reconciled by a versioned proof
crosswalk in the pinned files. We do not infer which
supersedes which, or that a proof cannot exist outside
the deposit. A reviewer should request the authoritative
companion revision and exact replacements for remaining
markers. Conversely, counting every outstanding string
as an independent necessary gap would be wrong: some
belong to optional routes, stronger conclusions or
statements explicitly not used by the active route.

\paragraph{The fifteen imported readings.}
Draft pp.142--143 labels them Pr1--Pr15. In order,
their subjects are:
\begin{enumerate}
\item large-field normalization/Wilson-term bounds,
CMP122-II (1.11), (1.23)--(1.48);
\item analyticity of old terms beside a new region, p.375;
\item the resummed boundary estimate, p.376;
\item locality and hull bookkeeping, pp.379, 390;
\item an estimate whose details are omitted, p.360;
\item interrupted integrations, CMP122-I p.191;
\item prepared-data admissibility, CMP122-I pp.177--179;
\item regularity on consecutive small-field differences,
CMP119 p.275;
\item transfer to a multiscale cover, CMP119 pp.276--277;
\item complex-domain inclusion, CMP109 (1.17);
\item a second regularity inequality, CMP109 p.279;
\item companion estimates asserted by analogy in the
propagator/averaging papers and CMP109 (3.26);
\item lowest-tower estimates, (3.35), (3.37) of
the draft's reference [5];
\item a background-field estimate derived from an
uncomputed display, reference [6], (189)--(190);
\item regularity of the unintegrated large-field core,
used for uniqueness.
\end{enumerate}
This is the draft's evaluation of its imports, not
our verdict that the originals are false or that
proofs cannot exist elsewhere. Domains and uniformities
must be compared in each application.

\paragraph{The uniqueness ledger.}
Draft pp.148--149 identifies three written statements
under independent verification, U1, U3 and U5;
two not yet written, U4 and U6; and a condition to
be added to a definition. These concern inherited-domain
extension, two-cutoff comparison of retained factors,
interior backgrounds, corrections near healed regions
and interrupted Gaussian factors. The draft says its
original U1 form was too strong and uses a corrected
form. Group-specific \(N\geq3\) work is separately
unperformed. These qualifications belong with any
citation of torus uniqueness.

\subsection{Navigation through the original mathematics}
\label{hrb-raw-map}

The raw account is not merely a classical calculation.
Its proposed chain contains full densities, large-field
histories, source estimates, continuum limits and a
witness. These locators lead to actual interfaces:
\begin{description}
\item[Conventions and quantifiers.]
I:11, 401 and 1050 begin Chapters 0, 1 and 2.
The action, physical pairing and left-exponential
expansion are at I:12756--12760, 15471--15478 and
15643--15648. The last expands at \(\exp(i\eta A)U_0\);
away from zero it is not the ordinary-potential
Hessian without a coordinate-current correction.
\item[Classical response and free currents.]
I:13681, \hrbpath{4.5:lem-free-current-operators};
I:14496, \hrbpath{4.5:thm-chart-existence};
I:19578, \hrbpath{4.5:thm-operator-bounds}.
The domains include graded spaces and independent
currents. Our homogeneous physical stationary
bounds do not verify all those domains.
\item[Averages and pivots.]
I:24240, \hrbpath{4.6:lem-averaging-locality};
I:24320--24325 gives corrected composition,
I:24349--24369 its flat differential.
I:46960 and following gives pivot/constraint
linear algebra. II:9032--9075 locates incident-bond
Dirichlet geometry, II:9779--9791 its constraint.
\item[Complete representation.]
I:61097, \hrbpath{4.20:thm-representation};
I:73545, \hrbpath{4.21:rem-global-chain},
specifies four global carriers.
I:92401--92620, \hrbpath{4.22:lem-current-localisation},
retains the complex-decoupling current correction.
\item[Sources, per-ball bounds and limits.]
II:2791, \hrbpath{5.1:thm-correlation-theorem};
II:60348, Chapter 6;
II:61205, \hrbpath{6.2:lem-chessboard-estimate};
II:64630, Chapter 7;
II:66412, \hrbpath{7.4:thm-infinite-volume}.
Its \(\mathbb R^4\) statement is part of the baseline;
it takes sourced estimates as inputs.
\item[Witness, uniqueness, rotations.]
II:68374 begins Chapter 8; III:5 begins Chapter 9.
III:25268, \hrbpath{9.2:thm-carry}, is an incomplete
uniqueness-path statement. III:75922 begins Chapter 10;
III:76512, \hrbpath{10.1:thm-rotational-identity}.
The far-flux assumption reappears at
III:110825--110835.
\end{description}
These are navigation points, not exhaustive local
proof checks. Shared words such as minimizer, covariance
and locality do not identify our results with the
source's full-history representation.

\subsection{Published results and transfer obligations}
\label{hrb-published}

Page numbers below are journal pages. Every cited theorem
retains its transformation, analytic class, current
variables, boundary conditions and coupling hypotheses.

\paragraph{Small fields and extraction.}
Ba{\l}aban, CMP109 (1987), 249--301,
\url{https://doi.org/10.1007/BF01215223},
Theorem 3, p.264, is the conditional small-field induction
for (1.1)--(1.22). Its separate Theorem 2, p.259,
announces the tuned \(SU(2)\) trajectory. Ward analysis
and full Wilson subtraction are in Sections 4--5:
(4.9)--(4.15), (4.21)--(4.45), (5.17)--(5.44).
Equation (1.19) requires gauge invariance of localized
coefficients. A total gauge identity in a different
activity representation is not literally that hypothesis.

CMP116 (1988), 1--22,
\url{https://doi.org/10.1007/BF01239022},
completes the small-field induction. Sections 1--2
localize and integrate the renormalized exponent;
(2.41) and pp.21--22 recover the inductive bound
and add the Gaussian normalization. This is a
substantive closure theorem under small effective
couplings, not just a formal loop expansion.
Our total-kernel extraction and quotient are
alternative arguments with separately stated inputs.

\paragraph{Complete densities.}
CMP119 (1988), 243--285,
\url{https://doi.org/10.1007/BF01217741},
defines full pushforward and the additional large-field
operation in (0.1)--(0.2); Section 2 specifies its
class; Theorem 1 is on p.262. The introductory
unnumbered statement is on p.245. Interfaces include
(2.18), histories; (2.19)--(2.22), large-field
operations; and (2.34)--(2.42), regularity/boundaries.
In particular (2.38), p.261, requires local gauge
potentials and first derivatives, not merely small
plaquettes. General averages satisfying its axioms
are allowed, so a composite map is not excluded by name.

CMP122-II (1989), 355--392,
\url{https://doi.org/10.1007/BF01238433},
Theorem 1, p.355, completes that class under a small
common ceiling for all effective couplings.
The concluding bounds are pp.390--391, (1.98)--(1.102).
Its introduction says the separate perturbative
trajectory proof was still unpublished in 1989 and
distinguishes later observable analysis.
Our search did not locate a later primary proof
of that precise theorem; this is a bounded search,
not proof of nonexistence. The original CMP122-I
formulas were not newly audited in full when its
primary PDF endpoint was inaccessible.

\paragraph{Universality and distinct constructions.}
A familiar coefficient cannot be inserted for an
unmatched determinant, cutoff or gauge normalization.
One-loop lattice comparison theorems, for example
Adams--Lee,
\url{https://arxiv.org/abs/0709.0781},
have hypotheses on their integrals. They do not
directly identify our constrained determinant.
Magnen--Rivasseau--S\'enoor's construction with
an infrared cutoff, CMP155 (1993), 325--383,
\url{https://doi.org/10.1007/BF02097397},
is a distinct regularized construction. Neither
citation is treated as a theorem about our exact
macrostep without identification.

\subsection{Our additions: objects, results and boundaries}
\label{hrb-additions}

Here addition means a derivation or interface produced
in this workspace, not a priority or novelty claim.
Gaussian integration, reflection positivity, forest
interpolation, cluster expansions, analytic implicit
functions, Ward identities and elliptic estimates
remain standard tools. The notebook manifest separates
our uses from the cited model-specific results.
The documents are research proof drafts with internal
audits, not accepted replacements for the literature.

\paragraph{01--05: continuum tools and explicit Gaussian interfaces.}
Notebook 01 gives a conditional compactness and
observable-reconstruction route, factorial-cumulant
issues and a stronger positive-energy test.
For example,
\(\langle v,e^{-tH}v\rangle\leq q\|v\|^2\), \(q<1\),
forces spectral weight at positive energy but
neither a gap nor an interacting gauge theory.
Notebook 02 develops same-limit diagonal reflection
positivity for pure Wilson theory using actual
slab gluing and comparison at fixed lattice coupling.
Its continuum application still needs common
observable limits and the relevant estimates.

Notebooks 03--05 isolate current identities,
Gaussian large-field tails, strict background
locality under local precision perturbations and
simple interaction cases. The forest expansion
differentiates Gaussian densities/covariance
parameters rather than sharp indicators without
justification. Mark localization and strict
background dependence are different results.
A product domain is stronger than an arbitrary
weak source-field ball. These modules do not
import every CMP119/CMP122 orbit domain.

\paragraph{06--17: a homogeneous classical and Gaussian base.}
The flat Maxwell work proves lower bounds,
compact right-inverse lifts and upper bounds for
the effective quadratic form \(q_M\), uniformly
in prior homogeneous refinement \(M\).
Alias formulas and pole cancellations give
exponential locality and a free refinement
limit with \(O(M^{-2})\) convergence in the
stated weighted norms. These are free-reference
limits, not interacting ultraviolet limits.

Subsequent modules identify the source-compatible
flat extension and gauge, construct fine inverses
and current responses, establish pointwise and
equal-exponent regularity, and control nonlinear
fixed-reference and moving-reference averages.
Ordinary-potential and left-exponential Hessians
are related with their current term retained.
A small global physical \(C^{1,\alpha}\)
background admits a local constrained Wilson
minimum, gauge-slice coverage and ordinary-constraint
transfer. Its classical action obeys \(DV=m\);
weighted response localizes its Hessian.
These statements do not cover arbitrary exterior
data, holes, mixed histories or independent
free-current domains.

Gaussian roughness is a negative interface
that must be retained. In four dimensions,
typical fine quantum fluctuations are outside
a fixed small physical \(C^{1,\alpha}\) ball;
the curl trace gives root-mean-square size
of order \(\eta^{-2}\). At a fixed current
step, correctly rescaled blocks instead
have uniform covariance bounds. Notebook 17
uses this normalization for joint tails.
Classical smooth-background estimates cannot
simply be applied to a full fine Gaussian sample.

\paragraph{18--22: local integration and strict classical support.}
Notebook 18 supplies an overlapping-interaction
Gaussian theorem with explicit support norms
and complex covariance control. Notebook 19
constructs the actual stationary-background
pivot chart and reduced precision, retaining
constraint curvature and coordinate-current
curvature. The independent retained-link
window is literal. Haar, ambient straightening
Jacobian, exponential-gauge factor and
background determinant are separately tracked.
Notebook 20 applies these to the bare local step.

Notebook 21 obtains strict support expansions
for the classical action through a simultaneous
field/multiplier fixed point, using smooth
physical cells, mixed insertion bounds and
the physical \(\eta^4\) pairing. It avoids
summing \(M^4\) coordinate coefficients and
does not assume a uniformly bounded ordinary
averaging adjoint where such a bound fails.
Notebook 22 transports inherited activities
through a stationary map with the explicit
multivariate Cauchy cost. Bounded substitution
does not automatically contract a norm or
preserve a radius through infinitely many steps.

\paragraph{23--28: covariance, re-entry and a common local domain.}
Notebook 23 treats actual gauge covariance
and ordinary representatives. Notebooks
24--26 treat inherited increments, pre-gauge
smooth representatives and auxiliary grouping.
The physical blocking factor and auxiliary
polymer grouping are distinct. The same
block labels are not counted twice as newly
compressed. Cardinality weights, fiber
dimensions, radius loss and thresholds
are budgeted explicitly.

Notebook 27 constructs the one-shot constraint
\(F_R=\mathcal M_\ell^s\), \(R=\ell^s\),
with a selected direct-\(R\) contour gauge,
actual pivots, measure factors and retained
window. Constraint composition does not
identify that window with successive
primitive windows. Notebook 28 proves a
common-domain local integration theorem
with joint background/fluctuation supports.
An inherited interaction is split into its
background value and fluctuation increment;
the former has no artificial coupling factor.
Equivariant measurable completion is useful
for an already defined global invariant
density. It does not turn a local analytic
germ into the Wilson density outside its domain.

\paragraph{29--31: extraction and finite local induction.}
Low-degree Ward cancellations are performed
on total kernels before localization.
Subtracting the full classical reference
preserves the nonlinear identity; subtracting
only its quadratic jet does not.
The total quadratic, cubic and quartic
remainders have the low-momentum orders
needed by the spatial Taylor argument.
Finite-torus versions keep wrapping defects
and the actual finite coefficient, rather
than imposing an infinite Fourier identity.

The activity quotient identifies normalized
totals, tracks vacuum constants, admits
controlled periodization/local limits and
supports quantum descent. The classical
semigroup identifies totals modulo the
allowed gauges, not canonical carrier arrays
or sequential quantum cutoffs.
Finite induction selects equivalent small
representatives with strict margins and
compensates scalar changes in the vacuum
factor. Its initial amplitude theorem was
conditional on a weak coupling sequence.
Completeness of an analytic interaction
space is not construction of a physical
probability law.

\paragraph{32--37: the actual determinant and positive local drift.}
These notebooks retain coarea normalization,
fine/target orbit densities, multiplier
curvature and actual stationary response.
Point-pinned ghosts, globally based quotients
and block mean-zero gauges are distinguished.
The moving-reference scalar construction yields
\[
 C-\bar C=D_cZ,\qquad \bar C D=D_cB,
 \qquad BB^T=K^{-4}I .
\]
This gives an exact gapped covariant
determinant pair at fixed \(K\), rather
than discarding a singular global-base
factor. Soft scalar/vector reductions
isolate coarse factors and bounded heat
ends. Local trace projections and ordered
Duhamel/resolvent products are retained:
a local trace does not have unrestricted
global cyclicity.

Notebook 37 claims, within the upstream
working chain, the actual coefficient estimate
\[
 \beta(\mathcal D_{M,K})
   =22\alpha_g\log K+O(1),\qquad
 \alpha_g=\frac{C_H}{48\pi^2},
\]
uniformly in prior homogeneous depth and
every applicable jointly obtainable
infinite structural system. Its Maxwell
coefficient is twice the normalized
affine \(t^2\) coefficient; \(C_H\) is
the adjoint charge-square sum in the
specified color convention. Numerical
comparison with another convention
requires conversion of action, trace
and extraction functional.

The comparison is direct: actual base
to Wilson has integrable \(O(s^{-1})\)
error, and fixed Wilson to continuum
has integrable \(O(s^{-1/2})\) error.
Joint compactness transfers a positive
block margin to actual finite coefficients.
Real finite-coupling continuity and
clipped shooting then produce arbitrarily
long finite weak homogeneous local
histories with a fixed weak endpoint
and bare coupling of order \(m^{-1/2}\).
This is a substantial local endpoint,
not a full-Wilson trajectory theorem,
a Schwinger-function limit or a mass gap.

\paragraph{38: exact statements about the full bare measure.}
Notebook 38 adds two different kinds
of result. A specified measurable coarse
map admits an exact positive good/remainder
decomposition retaining all histories
and normalization. Identification of
its good term with the analytic germ
requires literal containment of the real
integration window in the input domain.

Independently, direct compact-group
and positivity arguments control the
normalized bare Wilson measure.
Site-plane reflection positivity and
Schatten gluing give joint bad-plaquette
bounds. A fixed-box conditional Laplace
theorem uses an exact tree quotient,
a compact global energy gap and
cancellation of Gaussian volume factors
to obtain a normalized \(Ce^{-cT^2}\)
tail. A boundary-only gauge combines
these into joint bounds for separated
boxes with explicit local minimizer
windows. Constants are fixed-box
constants, not uniform in growing
components.

\paragraph{39: normalized local replacement in the bare Wilson measure.}
The final extension is
\hrbpath{outputs/ym4-normalized-box-replacement/main.tex}.
For an arbitrary finite family of
vertex-disjoint fixed-size boxes,
the replacement preserves the original
exterior marginal and raw partition
normalization. At a good boundary it
conditions the actual conditional law
on its minimizer window, dividing by
the exact success probability
\(1-p_i\geq1-a_g\); bad-boundary
conditionals are unchanged.

An exact positive failure expansion
and coupling control the joint local
law, including the exterior. For a
selected set of boxes \(A\), the
discrepancy is bounded by
\[
 \mathbb E\!\left[1-\prod_{i\in A}(1-p_i)\right]
 \leq |A|a_g ,
\]
independently of the total number of
replaced boxes. Discrepancy and bad-boundary
marks have the stated joint bound
\((a_g+\beta_g)^{|S|}\), and bounded
local sources/cumulants can be transferred.
These are actual bare-measure assertions.
They do not concern a nonlinear
RG constraint fiber, growing components
or an already renormalized density.
The final theorem's precise conditions
and internal audits remain part of
the handoff, rather than being erased
by this summary.

\subsection{Corrections and superseded routes}
\label{hrb-corrections}

Historical notes remain for traceability.
Their next-task and candidate language
is chronological. Later repairs and
counterexamples take precedence.
The following changes are particularly
important for a reviewer.

\begin{enumerate}
\item \textbf{No weakened baseline.}
The repository did claim \(\mathbb R^4\)
existence. Its definition already
excluded a two-point inference of
non-Gaussian interaction. Our additions
must not be presented against a
fabricated weaker starting claim.

\item \textbf{The mandatory three-map
source obstruction was withdrawn.}
An early reading of Raw II Section 5.6
treated three localized response maps
as the next unavoidable obstruction.
Raw I:73545,
\hrbpath{4.21:rem-global-chain},
uses four global carriers instead.
The correction is in
\hrbpath{work/source_stability_research.md},
opening and Section 9.
The conditional local estimates remain;
their asserted necessity does not.
The two descriptions still need an
authoritative use-by-use reconciliation.

\item \textbf{Sharp windows were kept
distinct.}
The retained-link cutoff and additional
pivot indicators cannot be conflated.
Primary CMP116 p.12 excludes selected
pivots from its retained bond set.
The nonlinear argument is \(\Phi_V(g\xi)\),
not generally \(g\Phi_V(\xi)\).
A bad link implies Euclidean cell norm
greater than the same threshold.
The separate good-window inclusion
with a square-root dimension factor
cannot improve the bad-event tail.

\item \textbf{Tails are not localized
logarithms.}
Joint Gaussian bad-block bounds alone
do not give connected log activities.
Forest/cluster and strict-support
arguments are additional work.
Grouping arbitrary large overlapping
supports into Gaussian blocks can
lose degree summability; no routine
extension claim substitutes for the
actual overlapping-interaction proof.

\item \textbf{Norm and radius losses
were real.}
A fixed cardinality reserve does not
pay all Taylor coefficients for free.
Transport uses the exact product
Cauchy cost, common smaller domains
and a geometry budget. Auxiliary
grouping/pre-gauge transport address
these losses rather than rename them.
When auxiliary vertices disappear,
connected-hull estimates use Steiner
length. Conversion to a terminal
metric tree spends a factor of at
most two in the exponent.

\item \textbf{Positive norm, positive
energy and a gap differ.}
A nonzero reflected vector can have
a zero-energy component. A time
correlation ratio forces positive
energy but no gap. A constrained
Gaussian gap \(cK^{-2}\) is not a
Yang--Mills Hamiltonian mass gap.
Independent random constants and
generalized-free examples remain
valid continuum-interface cautions.

\item \textbf{Pressure and free
boundaries did not identify the
same symmetric state.}
Surface free energy does not determine
local states, and translation averaging
need not preserve positivity at a
fixed mirror. The reflection proof
therefore uses actual gluing and
comparison with a common-limit
obligation, not unrelated states
chosen for individual tests.

\item \textbf{Gauges and Hessians
needed exact corrections.}
The source scalar projection is not
automatically projection onto
\(\ker Q_0\); a constrained biharmonic
construction was required.
Ordinary and left-exponential
Hessians differ by current curvature.
Constraint elimination adds multiplier
curvature. Dropping either changes
the determinant.

\item \textbf{Pinned and unpinned
infrared factors cannot be exchanged.}
Point pinning can give a scalar gap
of order \(K^{-4}\), not \(K^{-2}\),
and its long-time logarithm differs.
A globally based scalar factor has
a constant-holonomy anomaly with
forced vector compensation.
The later mean-zero pair is a
different exact decomposition, not
proof that the earlier factors
separately had the desired coefficient.

\item \textbf{A false difference-square
formula was replaced.}
Zero row sums alone do not give
the naive component-valued common
difference-square identity.
Component-skew/odd terms obstruct it.
The corrected proof uses ordered
Laurent factorization with zero
constant and first moments.
The amended covariant Hessian and
soft-vector/tail modules are operative.

\item \textbf{Separate uniform
relative vertices were not automatic.}
Pure-gauge and hierarchy-interface
examples give growing relative
vertices while full heat traces
remain invariant. Current arguments
use covariant densities, localized
forms and combined identities rather
than the rejected componentwise bound.

\item \textbf{Affine probes need a
total-kernel interpretation.}
A polynomial Cartan probe is not
a bounded periodic background.
Its use needs localized Taylor
kernels and polynomial moments.
A convenient local density on an
affine field need not equal the
total Maxwell moment. Reanchoring,
covariance and the limiting order
are indispensable.

\item \textbf{Infinite kernels are
not arbitrary finite kernels.}
Large-\(K\) magnetic/Bloch arguments
use a jointly chosen infinite
structural system or explicit
periodizations. Finite winding
errors and actual coefficient
transfer are separate.
Finite determinant algebra alone
does not justify thermodynamic
derivative interchange.

\item \textbf{Full-measure shortcuts
fail.}
Small plaquettes on a simply
connected box give a tree-gauge
link bound proportional to a
filling area. Alternating small
curvature rules out an automatic
mesh-uniform physical
\(C^{1,\alpha}\) chart.
Flat torus holonomies obstruct
an arbitrarily small global
potential. Raw link marginals
are Haar by gauge invariance,
so raw near-identity link
probability is the wrong event.
\end{enumerate}

\subsection{Remaining existence interfaces and audit priorities}
\label{hrb-remaining}

The local positive drift changes the
next research task; it does not
erase the remaining implications.

\paragraph{Full densities and
normalized large-field replacement.}
One needs a full match to the
CMP119/CMP122 class or a direct
replacement theorem for the chosen
composite map. It must handle
relevant bad components, exterior
data, holes, retained terms,
interrupted operations and ages,
with numerator/denominator
cancellation and summable bounds.
The weak couplings of our local
recursion are not automatically
the coefficients of that full
recursion. The averaging axioms
may permit the map, but all used
variational/domain hypotheses
still require verification.

The new normalized fixed-box
replacement is a useful concrete
piece of this problem. It does
not identify fixed-boundary
conditionals with conditionals
on nonlinear block averages,
or establish iteration on
post-RG densities. An exact
all-history identity is not
a uniform norm bound on
those histories.

\paragraph{Observables and a common
continuum limit.}
An external observable source
changes full densities,
counterterms and large-field
operations. Its cutoff/volume
uniform control is not implied
by holomorphy in a small
background connection.
Analytic activity compactness
is not tightness of a quantum
observable law. A common
hierarchy must retain consistency,
positivity, symmetries and the
growth required by the intended
reconstruction. Coincidence
extensions and the observable
algebra must be specified.

\paragraph{Identification and
infrared properties.}
A scalar action-density subsector
with a running two-point witness
is not alone identification of
an interacting Yang--Mills
theory. Clustering, vacuum
uniqueness and a mass gap are
separate infrared questions.
The source draft expressly
leaves them open; our local
ultraviolet work does not close
them. General color formulas
and bare \(SU(N)\) probability
lemmas do not extend every
upstream \(SU(2)\) input to
all compact simple groups.

\paragraph{Efficient audit order.}
First freeze the target theory
and authoritative companion
revision. Verify normalizations,
finite measure identities and
domains before long estimates.
For our series, inspect the
classical base, strict support
norms, common-domain quantum
theorem, finite extraction/
induction, determinant reduction/
comparison, then finite drift
transfer. Audit the full-measure
interface independently; the
length of the local chain
does not prove it. Historical
audit files guide checks but
are not independent certificates.

\subsection{Companion files and citation conventions}
\label{hrb-companions}

This report has a plain-text
counterpart,
\hrbpath{work/handoff_repository_baseline.txt},
and a machine-readable source map,
\hrbpath{work/handoff_repository_source_map.json}.
The current detailed obligation
ledger is
\hrbpath{work/final_push_existence_dependency_audit.txt}.
Primary source maps include
\hrbpath{work/renormalized_induction_primary_sources.md},
\hrbpath{work/full_measure_large_field_interface_research.txt},
and
\hrbpath{work/mean_zero_one_loop_source_interface.txt}.
The source-route correction is
in \hrbpath{work/source_stability_research.md}.
The notebook manifest identifies
authoritative assembled files
and the final extension's audits.

The repository and published
articles should be cited for
their own results and definitions.
Our claims should be cited as
research drafts with hypotheses.
This distinction remains necessary
even where a calculation seems
elementary or several internal
readers agree. The handoff supplies
a reproducible theorem-level map,
new explicit derivations and
remaining interfaces. It does
not declare the existence
problem solved.

\section{The local construction added in this project}
\label{hdc-local}

The main route is
\[
 \begin{gathered}
 \text{flat coercivity and classical minimizers}\\
 \longrightarrow\text{actual local quantum integral}\\
 \longrightarrow\text{Ward extraction and finite induction}\\
 \longrightarrow\text{actual positive determinant drift}\\
 \longrightarrow\text{arbitrarily long finite weak local histories}.
 \end{gathered}
\]
The catalogue supplies every intermediate notebook, rather than
replacing this chain by an appeal to perturbation theory.

\subsection{Flat and nonlinear classical analysis}

Notebooks 06--17 develop flat coercivity, constrained lifts,
weighted inverse kernels, refinement limits, pointwise regularity
and nonlinear response estimates for the homogeneous classical
minimization problem. The construction tracks physical pairings,
contour constraints, gauge directions, and the relation between
a prescribed coarse field and its fine minimizer.

Their endpoint is a family of local minimized classical actions
\(S_M\), analytic charts, and localized response kernels near
the trivial orbit. The support estimates matter because quantum
integration must preserve spatial summability, not merely
analyticity on a finite-dimensional ball.

These results concern regular homogeneous backgrounds. They do
not give a minimizer theorem on every determining set with
holes, arbitrary exterior boundary fields, or independent
complex currents. Their estimates cannot be used on those
domains by renaming variables.

\subsection{The actual local quantum integral}

Notebooks 18--28 formulate and control one fixed composite
ordinary coarse map \(F_R=\mathcal M_\ell^s\), its chosen
direct-\(R\) contour gauge, and its retained-coordinate
small-field sector. The nonlinear constraint chart
\(\Phi_V(z)\) keeps the retained coordinates unchanged and
solves the pivot variables. Its integration field is
\[
                 U=e^{i\Phi_V(g\xi)}V,
\]
not \(e^{ig\Phi_V(\xi)}V\). This distinction affects
higher-order interactions and the exact Jacobian.

The pullback includes Haar densities, the square ambient
straightening, the pivot constraint derivative, the contour
gauge penalty, the actual reduced Hessian, and the independent
retained link/color-ball window. A rectangular determinant of
the retained embedding cannot replace the constraint Jacobian.

The compact gauge insertion uses the exact normalizer
\[
 c(g)=\int_G\exp[-(1-\operatorname{Re}\tau h)/g^2]\,dh.
\]
For a literal represented input density the local output
has the following factorization, with notation abbreviated
but the normalization displayed:
\begin{align*}
 &Z_{\rm in}\,c(g)^{-n_s}h_0^{n_e-n_c}D_R(V)\,
 g^{n_z}(2\pi)^{n_z/2}\det A_M(V)^{-1/2}\\
 &\quad{}\times e^{-S_M(V)/g^2+E_{\rm raw}(V)}
 \mathbb E_{A_M(V)^{-1}}
       [\chi_{\rm window}\chi_{\rm extra}e^{I(V,g,\xi)}].
\end{align*}
Here \(n_z=(\dim\mathfrak g)(n_e-n_c)\), \(D_R\) is the
actual pivot factor, \(h_0\) is the Haar density at the origin
in the chosen Lie coordinates, and \(I\) contains all higher
action, Haar, straightening and inherited fluctuation terms.
Individual notebooks define every factor without shorthand.

The common-domain argument prevents uncontrolled shrinkage of
the analytic field ball at every step. It uses the actual
pre-gauge representative and transports the \emph{entire}
inherited insertion, including its background value.
It does not merely change a fluctuation variable and forget
its cutoff. Strict connected support families and
overlapping-carrier expansions provide the spatial norms.

Composition of primitive constraints identifies the ordinary
coarse map. It does not identify one composite retained window
with an intersection of sequential primitive windows. Also a
holomorphic continuation of a local integral is not a
representation of the full positive measure outside its
real integration domain.

\subsection{Ward extraction and finite local induction}

Notebooks 29--31 use Ward identities of the total functional
to isolate its marginal Maxwell contribution. They do not
assume that each arbitrarily selected activity is separately
gauge invariant. The remaining interactions are relocalized
and controlled in a suitable representative of the exact
finite total.

That representative need not be the original canonical cluster
array. Equivalence of the total and its compensating vacuum
constant are tracked: changing \(E\) to \(E+c\) requires
changing \(Z\) to \(Ze^{-c}\). Dropping that factor would
destroy a later full-measure normalization comparison.

The finite induction bounds the local remainder for permitted
weak coupling sequences and sufficiently large compatible
periods. It does not itself force the couplings to remain
weak. Positive drift and finite shooting supply that part.

\subsection{The actual logarithmic drift}

Notebooks 32--37 compute the leading coefficient for the
actual constrained determinant. An exact soft reduction
expresses it as a fine scalar/vector pair and coarse factors.
The coarse Maxwell contributions and both heat-integral ends
are bounded uniformly in preceding depth.

On the remaining band \(1<s<K^2\), source-localized resolvents
remove soft penalties with an integrable error. Nonlinear
Ward jets through field degree four compare the actual
effective-action vector base with the Wilson base;
left-relative chart-current terms remain present.
A direct Wilson-to-continuum comparison controls the
remaining error and fixes signs and color multiplicities.

With \(\tau(H^2)=1\) and the full adjoint charge trace,
\[
 C_H=\sum_{q\in\operatorname{spec}(\operatorname{ad}_H)}q^2,
 \qquad \alpha_g=\frac{C_H}{48\pi^2},
\]
the working theorem is
\[
       \beta(\mathcal D_{M,K})=22\alpha_g\log K+O(1).
\]
The error is uniform in old depth and every applicable
joint infinite structural system, under the classical and
local quantum inputs explicitly used in the series.
For \(\tau=\operatorname{Tr}/N\), the color calculation
gives \(C_H=2N^2\); this does not establish earlier
construction inputs for every \(SU(N)\).

Here \(\beta=2[t^2]\) on the normalized affine probe.
The signed contributions of the mean-zero scalar and vector
factors to \(\mathcal D\) have coefficients
\(2\alpha_g\) and \(20\alpha_g\). An older point-pinned
ghost decomposition distributes logarithms differently;
its scalar coefficient cannot be combined with the new
mean-zero vector factor.

\subsection{Finite weak histories and the recursion actually tuned}

Finite moment transfer converts the positive margin in every
joint structural limit to a uniform finite block margin.
The geometric ratios are fixed first; a block of \(q\)
steps then dominates the bounded comparison error; periods
and the weak ceiling are chosen so the actual errors do
not consume that margin.

A continuous clipped finite recursion is shot to a fixed
weak endpoint \(\Gamma_*\), where \(\Gamma=g^{-2}\).
Positive future blocks show clipping did not occur on
the selected trajectory. With \(m\) blocks,
\[
 (\Gamma_*+mqD_{\rm abs})^{-1/2}
 \le g_{\rm bare}\le
 (\Gamma_*+m\delta/4)^{-1/2}.
\]
This supplies arbitrarily long finite \emph{homogeneous local}
weak histories. It does not claim uniqueness, consistency
between horizons, or that a full large-field density has
the same extracted recursion. That identification remains open.

\section{The full-measure additions}
\label{hdc-fullmeasure}

Notebooks 38--39 add exact finite positive-measure results.
Their bare Wilson probability theorems are independent of
the earlier determinant work. The local-germ embedding uses
the already established coarse chart and Ward identities.

\subsection{Exact positive history operators}

Fix a measurable equivariant completion of the actual local
coarse map, using its original mean on its good set and
its original reference string outside. This specifies a
global coarse observable of the original fine measure;
it does not invent a global analytic effective action.

The compact based-gauge insertion gives a probability
kernel preserving the coarse field. An actual retained
sheet yields a local density with
\[
                  F_*\mu\ge 1_D\ell_\mu\,m_C.
\]
For a nonnegative cutoff \(\theta\) supported in \(D\),
let \(m_\theta(W)=\int\theta(W^v)\,dv\). Define
\[
 \mathcal A_g\mu(dW)=
 \frac{\int\theta(W^v)\ell_\mu(W^v)\,dv}{m_\theta(W)}\,m_C(dW),
\]
with value zero when \(m_\theta=0\). The invariant
absolutely continuous density of the true pushforward
bounds every numerator orbit value. Hence
\[
       0\le\mathcal A_g\mu\le F_*\mu,\qquad
       \mathcal R_g=F_*-\mathcal A_g\ge0.
\]
These linear invariant operators retain singular parts
in the remainder. Their finite products account for
every history, raw mass and observable numerator exactly.

Local-germ equality uses a finite orbit-overlap lemma
and the literal invariant-ball cutoff Ward identity.
Its radius may shrink with graph diameter. Every
reconstructed real field must also lie where the input
measure actually has the claimed density.
Holomorphic continuation alone supplies neither condition.

\subsection{Actual Wilson rarity and conditional normalization}

For site-plane cells of fixed side \(b\ge3\) and torus
periods \(n_i=2m_i b\), exact Haar slab gluing and even
Schatten H\"older give a chessboard inequality.
A lower bound on the \emph{full} partition function yields
\begin{align*}
 \mu_g\left(\bigcap_{j\in S}B_j(\delta)\right)
 &\le z_b(g,\delta)^{|S|},\\
 z_b&=\min\{1,e^{-\delta^2/g^2+
       b^d(d\dim G\log(1/g)+C_{G,d})}\}.
\end{align*}
Each \(B_j\) forces an interior plaquette energy
\(1-\operatorname{Re}\tau U_p\ge\delta^2\).
Finite coloring gives prescribed-plaquette bounds with
exponent loss \(Q_b=b^d\binom d2\).
This is a normalized interacting Wilson estimate,
uniform in surrounding volume.

For a fixed box \([0,R]^4\) with small prescribed
boundary edges, a forest rooted at every boundary
vertex removes the interior gauge redundancy exactly.
At identity boundary its representative has a unique
zero-action field and a positive Hessian. Compactness
and the implicit-function theorem give a unique global
conditional minimum for nearby data. Numerator and
denominator have the same local factor \(g^n\);
the outer compact region has a fixed energy gap. Thus
\[
      \mu_{V,g}(W_{V,g,T}^c)\le Ce^{-cT^2},
      \qquad 1\le T\le r/g.
\]
The window uses an adjoint-invariant edge-max norm
about that conditional minimizer.

The boundary complex is \(S^3\). Boundary tree loops
have finite boundary-plaquette fillings, so small
boundary curvature supplies a boundary-only small
gauge. Conditioning on the exterior of vertex-disjoint
boxes makes their interiors independent. Combining
the estimates gives
\[
 \mu_g\left(\bigcap_{i=1}^m D_i\right)
 \le \left[Ce^{-cT^2}
       +24R^3 z_b(g,\varepsilon_R)^{1/Q_b}\right]^m.
\]
Here \(D_i\) is a bad boundary or a failed local
gauge/minimizer window. It is an event in the complete
bare measure. It is not yet the nonlinear averaged
constraint window of the iterated map.

\subsection{The final normalized replacement}

Notebook 39 conditions on the common exterior \(\omega\).
The interior laws factor as \(\prod_iQ_i^\omega\).
On a good boundary put
\[
 p_i(\omega)=Q_i^\omega(W_i^c)\le a_g,\qquad
 P_i^\omega=Q_i^\omega(\,\cdot\mid W_i).
\]
On a bad boundary take \(P_i^\omega=Q_i^\omega\) and
\(p_i=0\). The operation does not remove bad boundaries.

For the original exterior marginal \(\rho\), set
\[
          \nu(d\omega,dx)=\rho(d\omega)\prod_iP_i^\omega(dx_i).
\]
This preserves probability and, in the unnormalized
version, the original raw partition mass.
Its Radon--Nikodym factor is a product of good-window
indicators divided by their actual probabilities
\(1-p_i(\omega)\ge1-a_g\).

An explicit coupling changes an interior only on a
conditionally independent failure of probability \(p_i\).
For the exterior jointly with a selected set \(A\),
total variation, with convention
\(\sup_E|\mu(E)-\nu(E)|\), is exactly
\[
 \mathbb E_\rho\left[1-\prod_{i\in A}(1-p_i)\right]
 \le |A|a_g.
\]
The estimate is independent of the number of other
boxes. The full theorem gives sparse discrepancy
bounds, retains correlated bad boundaries, and
transfers bounded local sources with the same
normalized laws.

This completes a finite localized replacement theorem.
It does not cover arbitrary growing components,
repeated scales, or nonlinear coarse-average fibers.
Preserving normalization does not establish that the
homogeneous coupling recursion remains unchanged.

\section{The remaining theorem program}
\label{hdc-open}

\subsection{H2 Complete-density transfer for the chosen map}

The first major unresolved theorem is propagation of
a complete density for the same full Wilson measure
and coupling recursion as the local construction.
There are two routes.

An import route verifies the entire induction class
of Ba{\l}aban's complete-density theorem for the chosen
coarse map and operations. That theorem is conditional
on its own effective couplings staying weak.
Matching their expected size or one-loop sign is not
a match of the recursion. The primary-source map
appears in the baseline assessment and Notebook 38.

A direct route builds a spatial determining-set
decomposition with collars, a regular constrained
background, rough fluctuations and exceptional
regions. It proves exact good-component integration
with the actual window and Haar, gauge, pivot and
vacuum factors. It then bounds normalized replacement
on every allowed complementary component, with
constants summable in geometry and scale age.
Notebook 39 supplies a restricted fixed-box example.

The proof should state admissible exterior fields,
transition or holonomy data, constraints, conditional
measures, analytic domains and the component norm.
Its denominator lower bound must be compatible with
the numerator's boundary energy. An unrelated
Gaussian normalization cannot be inserted after
conditioning.

Several shortcuts are excluded. Unfixed raw link
marginals are Haar. Flat periodic fields can have
nontrivial holonomy. Even plaquettes of size
\(O(\eta^2)\) can have alternating physical curvature
whose adjacent norms prohibit every uniformly
bounded \(C^{1,\alpha}\) gauge. Small curvature
alone is not a smooth-background construction.

\subsection{H3 Observable and source estimates}

After complete-density control, external source
estimates are needed for a specified gauge-invariant
observable algebra. Analyticity in a connection
background is not analyticity in an external
composite-observable source. The latter changes the
density, local counterterms, translation properties
and possibly its large-field estimates.

A useful first target is renormalized action density
on reflected separated test supports. Its source
neighborhood must be uniform in mesh and volume
in the exact cells used by reflection. Mixed
derivatives must obey all-order distributional
bounds, for example factorial bounds in stated
test-function seminorms. Fixed-point-number bounds
with no controlled growth in that number do not
supply field reconstruction estimates.

Numerator and denominator must belong to the same
positive regulator construction. Notebook 39's
bounded-source theorem does not control renormalized
fields whose supremum norms grow with refinement.

\subsection{H4 A common continuum hierarchy}

Once actual model bounds are available, extract one
diagonal subsequence for a countable dense test
algebra and extend by uniform continuity. Products,
mixed insertions, subtraction conventions and
required contact extensions must be compatible
on that same subsequence. Compactness of analytic
interactions on a background ball is a different
problem from tightness or distributional
compactness of physical correlations.

Separated correlations may not determine every
distribution needed for local fields. A route
avoiding particular coincidence extensions should
state why; otherwise those extensions and
counterterms must be constructed compatibly with
positivity and covariance.

\subsection{H5 Positivity, rotations and reconstruction}

Reflection positivity must descend from the same
positive finite measures and limiting observable
algebra. Estimates for a substitute measure or
unrelated subsequence cannot simply be combined
with it.

Hypercubic symmetry is not \(O(4)\) covariance.
The reflection-based restoration mechanism has
conditional geometric and source-cell reductions
in the early notebooks. Its uniform model
estimates still need to hold for all transformed
supports, cell positions and all-order constants
used by that argument.

Osterwalder--Schrader reconstruction requires the
chosen hierarchy's precise regularity and
positivity hypotheses and an appropriate vacuum
sector. Unique vacuum or clustering needs its
own estimate. If existence can use a suitable
sector or decomposition instead, that choice
should be explicit.

\subsection{H6 Nontrivial Yang--Mills identification}

A positive beta coefficient does not establish
nonzero positive-energy states in the reconstructed
theory. A quantitative lower witness must survive
the same limit. In a reconstructed theory, a
strict inequality
\[
 \|\psi\|^2-\langle\psi,e^{-tH}\psi\rangle>0,\qquad t>0,
\]
would certify positive-energy content. The task is
to prove such a witness from the cutoff theory.

A non-Gaussian scalar composite is not by itself
an interaction witness: a quadratic composite of
a free Gaussian field is already non-Gaussian.
The algebra, tensor laws and short-distance
normalizations must be sufficient to identify
the limit as Yang--Mills. A sufficient
identification criterion should be stated and
proved; no speculative criterion is promoted to
a theorem in this dossier.

\subsection{H7 Physical mass gap and group scope}

The gap in an auxiliary constrained Gaussian,
typically at the blocking scale, is not a
positive gap in the interacting Hamiltonian
above the vacuum. A physical infrared estimate
must be uniform in the same continuum sequence.
An \(SU(2)\) construction does not settle the
all-compact-simple-group target. These are
separate from finishing the ultraviolet program.

\section{How a successor can use the bundle}
\label{hdc-continuation}

\subsection{A recommended first pass}

Read the start document, the baseline assessment
and catalogue entries 31--39. Then inspect
Notebook 37's actual coefficient theorem and
trace its upstream inputs, or independently
begin with Notebooks 38--39's bare Wilson
arguments. These routes need different
expertise and can be reviewed in parallel.

For the drift, prioritize relative
Hessian/current terms, nonlinear Ward jets,
the determinant decomposition, infrared and
source-localized traces, and the transfer from
joint structural limits to finite drift.
For complete-density transfer, prioritize
domains, coarse-map completion, retained-window
equality, exterior normalization and coupling
extraction.

\subsection{A concrete next research deliverable}

The next useful deliverable is a theorem for a
collared component with the actual averaged
constraint. It should specify:
\begin{enumerate}
\item the component, exterior data,
gauge/transition variables and compact
conditional integral;
\item the constrained background and retained
map with its full Jacobian;
\item an exact normalized replacement or
good/remainder identity preserving the fine
measure when components are summed;
\item a bound in the complete-density
induction's actual norm;
\item compatibility with neighboring regions,
scale changes, sources and coupling extraction.
\end{enumerate}

Verifying the published induction class may
be more efficient than rebuilding its
geometry. Choose by an exact hypothesis
comparison. A direct construction changing
the coarse map must redo its local
identification and drift transfer.

\subsection{Review and reproducibility}

Internal audits and successful compilation
are recorded, not called external
certification. A successor should look for
hidden domain enlargement, dependence on a
supposedly uniform parameter, an unrelated
conditional denominator, or identities
combined from different limits. These are
specific failure modes encountered and
repaired during this project.

Portable sources contain the assembled
notebooks and historical fragments.
Reproduction instructions give compile
commands; the manifest records delivered
bytes. Third-party articles are linked
rather than republished. The repository
can be fetched at its pinned commit and
checked against the recorded tree.

The outcome is a detailed foundation and
continuation program. The ambition remains
existence; the bundle makes the difference
between that ambition and the established
arguments inspectable.

\end{document}
